\subsection{Composition, Restriction, and Inverse Functions}

\begin{definition} \hlt{Composition, Restriction, Extension, and Identity}\\
Let $f:X\to Y$.
\begin{enumerate}[label=(\roman*)]
\setlength{\itemsep}{0pt}
\item for $A\subseteq X$, $f|_A:A\to Y$ is the \hlt{restriction} of $f$ to $A$;
\item if $X\subseteq U$, an \hlt{extension} of $f$ to $U$ is $\widetilde f:U\to Y$ with $\widetilde f|_X=f$;
\item if $g:C\to Z$ and $f(X)\subseteq C$, then $g\circ f:X\to Z$ is given by $(g\circ f)(x)=g(f(x))$;
\item $\operatorname{id}_X:X\to X$ denotes the \hlt{identity function}.
\end{enumerate}
\end{definition}

\begin{theorem} \label{thm:bijective-iff-invertible-function} \hlt{Bijective Means Invertible}\\
A function $f:X\to Y$ is bijective iff there exists a function $g:Y\to X$ such that $g\circ f=\operatorname{id}_X$, $f\circ g=\operatorname{id}_Y$.\\
When such a function exists, it is unique and is denoted $f^{-1}:Y\to X$.
\begin{proof}
\begin{enumerate}[label=\roman*.]
\setlength{\itemsep}{0pt}
\item[$(\Rightarrow)$] If $f$ is bijective, define $g(y)$ as the unique $x$ with $f(x)=y$. Then $g\circ f=\operatorname{id}_X$ and $f\circ g=\operatorname{id}_Y$.
\item[$(\Leftarrow)$] If such a $g$ exists, applying $g$ to $f(x_1)=f(x_2)$ gives $x_1=x_2$, and $y=f(g(y))$ gives surjectivity.
\end{enumerate}
If $g$ and $h$ are two-sided inverses, then $g(y)=g(f(h(y)))=h(y)$, so the inverse is unique.\\
\hlt{Structural Commentary}: A true inverse undoes $f$ on both sides.\\
\hlt{Key Insight}: Bijective means exactly one preimage per output.
\end{proof}
\end{theorem}

\begin{proposition} \label{prop:composition-restriction-calculus} \hlt{Composition and Restriction Calculus}\\
Let $f:X\to Y$, $g:Y\to Z$, and $h:Z\to W$.
\begin{enumerate}[label=(\roman*)]
\setlength{\itemsep}{0pt}
\item $h\circ(g\circ f)=(h\circ g)\circ f$.
\item $\operatorname{id}_Y\circ f=f$ and $f\circ\operatorname{id}_X=f$.
\item For every $A\subseteq X$, $(g\circ f)(A)=g(f(A))$.
\item For every $C\subseteq Z$, $(g\circ f)^{-1}(C)=f^{-1}(g^{-1}(C))$.
\item For every $A\subseteq X$, $(g\circ f)|_A=g\circ(f|_A)$.
\item If $f$ and $g$ are bijective, then $g\circ f$ is bijective and $(g\circ f)^{-1}=f^{-1}\circ g^{-1}$.
\end{enumerate}
\begin{proof}
\begin{enumerate}[label=(\roman*)]
\setlength{\itemsep}{0pt}
\item For $x\in X$, $(h\circ(g\circ f))(x)=h(g(f(x)))=((h\circ g)\circ f)(x)$.
\item For $x\in X$, $(\operatorname{id}_Y\circ f)(x)=f(x)$ and $(f\circ\operatorname{id}_X)(x)=f(x)$.
\item For $z\in Z$, $z\in(g\circ f)(A)\Longleftrightarrow\exists x\in A,\ z=g(f(x))\Longleftrightarrow z\in g(f(A))$.
\item For $x\in X$, $x\in(g\circ f)^{-1}(C)\Longleftrightarrow g(f(x))\in C\Longleftrightarrow f(x)\in g^{-1}(C)\Longleftrightarrow x\in f^{-1}(g^{-1}(C))$.
\item For $x\in A$, $((g\circ f)|_A)(x)=g(f(x))=(g\circ(f|_A))(x)$.
\item If $f$ and $g$ are bijective, then $f^{-1}:Y\to X$ and $g^{-1}:Z\to Y$ exist. For $x\in X$, $(f^{-1}\circ g^{-1})((g\circ f)(x))=x$; for $z\in Z$, $(g\circ f)((f^{-1}\circ g^{-1})(z))=z$. Thus $f^{-1}\circ g^{-1}$ is the inverse of $g\circ f$.
\end{enumerate}
\hlt{Structural Commentary}: Composition is pointwise; inverses reverse order.\\
\hlt{Key Insight}: Evaluate at an element.
\end{proof}
\end{proposition}
